%%%%%%%%%%%%%%%%%%%%%%%%%%%%%%%%%%%%%%%%%%%%%%%%%%%%%%%%%%%%%%%%%%%%%%%%%%%
%
% Generic template for TFC/TFM/TFG/Tesis
%
% $Id: desarrollo.tex,v 1.4 2015/06/05 00:05:18 macias Exp $
%
% By:
%  + Javier Macías-Guarasa.
%    Departamento de Electrónica
%    Universidad de Alcalá
%  + Roberto Barra-Chicote.
%    Departamento de Ingeniería Electrónica
%    Universidad Politécnica de Madrid
% 
% Based on original sources by Roberto Barra, Manuel Ocaña, Jesús Nuevo, Pedro Revenga, Fernando Herránz and Noelia Hernández. Thanks a lot to all of them, and to the many anonymous contributors found (thanks to google) that provided help in setting all this up.
%
% See also the additionalContributors.txt file to check the name of additional contributors to this work.
%
% If you think you can add pieces of relevant/useful examples, improvements, please contact us at (macias@depeca.uah.es)
%
% You can freely use this template and please contribute with comments or suggestions!!!
%
%%%%%%%%%%%%%%%%%%%%%%%%%%%%%%%%%%%%%%%%%%%%%%%%%%%%%%%%%%%%%%%%%%%%%%%%%%%

\chapter{Diseño e implementación}
\label{cha:design}

Este capítulo aborda el diseño e implementación de las técnicas metaheurísticas descritas para la resolución del Problema
del Viajante. El objetivo es trasladar los conceptos estudiados a implementaciones funcionales que permitan comparar el 
comportamiento de los diferentes algoritmos bajo unas condiciones comunes.

\section{Entorno de desarrollo}
\label{sec:dev_env}

El desarrollo e implementación de los algoritmos se va a realizar mediante el lenguaje de programación \textbf{\textit{Python}}.
Se ha escogido este lenguaje debido a su amplio uso en el ámbito de la computación científica y la optimización, así como a la 
disponibilidad de librerías que facilitan el tratamiento de datos y la implementación de los algoritmos estudiados.

Los experimentos se van a ejecutar en un ordenador personal, cuyas principales características de \textit{hardware} y
\textit{software} se incluyen en la tabla~\ref{tab:hw_sw}. La especificación del entorno permite establecer las condiciones
bajo las que se han obtenido los resultados experimentales y facilita la reproducibilidad de las pruebas realizadas.

\begin{table}[ht]
    \renewcommand{\arraystretch}{1.3}
    \caption{Características y elementos \textit{hardware} y \textit{software} del entorno experimental.}
    \label{tab:hw_sw}
    \begin{center}
        \begin{tabular}{|c|c|}
            \hline
            \textbf{Componente} & \textbf{Especificación}\\
            \hline
            Fabricante & HP \\
            \hline
            Procesador & Intel® Core™ i7-1065G7 CPU @ 1.30GHz  \\
            \hline
            GPU  & Intel® Iris® Plus Graphics \\
            \hline
            RAM & 16GB \\
            \hline
            Almacenamiento & 1TB \\
            \hline
            Sistema operativo & Microsoft Windows 11 Home (64 bits)   \\
            \hline
            Lenguaje de programación & Python 3.13 \\
            \hline
            Entorno de desarrollo & \textit{Visual Studio Code} \\
            \hline
        \end{tabular}
    \end{center}
\end{table}

\section{TSPLIB}
\label{sec:tsplib}

Para evaluar el comportamiento de los algoritmos es necesario disponer de un conjunto de instancias de referencia, que permita
realizar experimentos reproducibles y comparar los resultados obtenidos con otros trabajos. Para este propósito, se ha utilizado
\textbf{TSPLIB}, una colección de instancias de problemas de optimización combinatoria ampliamente utilizada en el estudio del 
Problema del Viajante, así como otros problemas relacionados \cite{tsplib}. Fue publicada en 1991, por el alemán Gerhard Reinelt
\cite{reinelt1991} \cite{wikipedia2026}.

La colección de instancias para el Problema del Viajante es accesible en \cite{tsp_dataset}. En ella se encuentra una lista de 
instancias estandarizadas ordenadas por orden alfabético. El nombre para las instancias consta de una palabra y un número:
\begin{itemize}
  \item La \textbf{palabra} representa, habitualmente, la inicial del autor que diseñó el problema, un acrónimo de la 
  procedencia de los datos o el tipo de problema real del que se extrajo.
  \item El \textbf{número} indica el total de nodos o ciudades que componen la instancia.
\end{itemize}
Por ejemplo, en el caso de la instancia \texttt{berlin52}, se tiene que dicha instancia cuenta con 52 ubicaciones reales de la 
ciudad de Berlín.

Si se descarga una de estas instancias, por ejemplo \texttt{berlin52}, se obtiene un fichero \texttt{.gz}, que al descomprimirse,
se obtiene el fichero en su formato \texttt{.tsp}. Este tipo de ficheros se puede abrir con cualquier editor de texto, ya sea el 
bloc de notas o un entorno de desarrollo cualquiera.

El fichero \texttt{berlin52.tsp} presenta el siguiente aspecto:
\begin{figure}[htbp]
  \centering
\begin{verbatim}
NAME: berlin52
TYPE: TSP
COMMENT: 52 locations in Berlin (Groetschel)
DIMENSION: 52
EDGE_WEIGHT_TYPE: EUC_2D
NODE_COORD_SECTION
1 565.0 575.0
2 25.0 185.0
3 345.0 750.0
4 945.0 685.0
5 845.0 655.0
6 880.0 660.0
7 25.0 230.0
...
50 595.0 360.0
51 1340.0 725.0
52 1740.0 245.0
EOF
\end{verbatim}
  \caption{Contenido de la instancia \texttt{berlin52.tsp}.}
  \label{fig:inst_info}
\end{figure}

Como se puede observar, cada instancia cuenta con los siguientes campos:
\begin{itemize}
  \item \textbf{NAME}: indica el nombre de la instancia, idéntico al del fichero.
  \item \textbf{TYPE}: el tipo de problema de optimización al que se corresponde la instancia. En todos los casos de este 
  trabajo será TSP.
  \item \textbf{COMMENT}: incluye una pequeña frase opcional informando sobre la instancia. En este caso, nos confirma que las
  52 localizaciones de la instancia pertenecen a la capital alemana.
  \item \textbf{DIMENSION}: incluye la dimensión de la instancia, esto es, el número de nodos.
  \item \textbf{EDGE\_WEIGHT\_TYPE}: indica cómo se determinan las distancias. En este caso, los nodos se presentan como \textbf{
  coordenadas en un plano euclideano de dos dimensiones} (\texttt{EUC\_2D}). Otros valores posibles son \texttt{EXPLICIT} o 
  \texttt{GEO}, entre otros.
  \item \textbf{NODE\_COORD\_SECTION}: es la información más importante de la instancia, puesto que se trata de las \textbf{
  coordenadas de cada una de las localizaciones}. En otras instancias, se incluyen los \textbf{pesos de las aristas} en su lugar.
\end{itemize}

Además de esta colección de instancias estandarizadas, TSPLIB también cuenta con una sección con los \textbf{resultados óptimos}
para cada una de estas instancias. Se puede consultar en \cite{tsp_best}. La existencia de esta sección es crucial, puesto que 
el conocimiento de la solución óptima es clave para calcular la proximidad de los óptimos propuestos por los algoritmos con el 
óptimo global, y así establecer unas métricas de eficiencia.

\section{TSPLIB95}
\label{sec:tsplib95}

Para conectar el uso del lenguaje Python y la colección de instancias TSPLIB, se va a utilizar una librería de Python llamada
\textbf{tsplib95}, desarrollada por Robert Grant. Se puede consultar su documentación en \cite{grant_tsplib95} y su código 
fuente en \cite{grant_github}, al ser una librería \textit{open source}.

Para instalar la librería, basta con ejecutar el siguiente comando en la terminal del \ac{IDE}:
\begin{verbatim}
  pip install tsplib95
\end{verbatim}
Una vez ejecutado, se puede consultar la versión de la librería, así como más información acerca de ella, con el comando:
\begin{verbatim}
  pip show tsplib95
\end{verbatim}
El resultado es el siguiente:
\begin{verbatim}
  Name: tsplib95
  Version: 0.7.1
  Summary: TSPLIB95 works with TSPLIB95 files.
  Home-page: https://github.com/rhgrant10/tsplib95
  Author: Robert Grant
  Author-email: rhgrant10@gmail.com
  License: Apache Software License 2.0
  Location: C:\Users\<usuario>\AppData\Local\Programs\Python\Python313\Lib\
    site-packages
  Requires: Click, Deprecated, networkx, tabulate
  Required-by:
\end{verbatim}
Se puede confirmar, además de la versión del \textit{software}, su autoría, licencia, y lugar de instalación en el sistema.

\subsection{Uso de la librería}
\label{subsec:lib_usage}

Para comenzar con el uso de tsplib95, se debe cargar una primera instancia. Esto se realiza mediante la función \texttt{load()}.
Por ejemplo, para cargar la instancia \texttt{berlin52}, se debe ejecutar:

\begin{verbatim}
  import tsplib95
  problem = tsplib95.load("berlin52.tsp")
\end{verbatim}

Asumiendo que el fichero \texttt{berlin52.tsp} se encuentre en el mismo directorio que el \textit{script} ejecutor de Python.

Una vez cargada la instancia, se pueden realizar una serie de operaciones para conocer más información de la instancia dentro 
del entorno de desarrollo. Se puede cargar toda la información del problema mediante la orden:

\begin{verbatim}
  print(problem.render())
\end{verbatim}

Se obtiene la misma información de la figura~\ref{fig:inst_info}. Se pueden extraer algunos de estos campos de manera individual,
como el nombre de la instancia:

\begin{verbatim}
  >>> print(problem.name)
  berlin52
\end{verbatim}

El formato en el que se presentan las distancias:

\begin{verbatim}
>>> print(problem.edge_weight_type)  
EUC_2D
\end{verbatim}

Una lista de los nodos de la instancia:

\begin{verbatim}
  >>> print(list(problem.get_nodes()))
  [1, 2, 3, 4, 5, 6, 7, 8, 9, 10, 11, 12, 13, 14, 15, 16, 17, 18, 19, 20, 
  21, 22, 23, 24, 25, 26, 27, 28, 29, 30, 31, 32, 33, 34, 35, 36, 37, 38, 
  39, 40, 41, 42, 43, 44, 45, 46, 47, 48, 49, 50, 51, 52]
\end{verbatim}

Se puede crear un objeto del tipo \texttt{networkx.Graph} a partir de la instancia:

\begin{verbatim}
  graph = problem.get_graph()
\end{verbatim}

A partir de este objeto, se puede obtener información de tipo grafo. Por ejemplo, información general sobre el propio grafo:

\begin{verbatim}
  >>> graph.graph
  {'name': 'berlin52', 'comment': '52 locations in Berlin (Groetschel)', 
  'type': 'TSP', 'dimension': 52, 'capacity': 0} 
\end{verbatim}

Su lista de nodos:

\begin{verbatim}
  >>> graph.nodes
  NodeView((1, 2, 3, 4, 5, 6, 7, 8, 9, 10, 11, 12, 13, 14, 15, 16, 17, 18,
  19, 20, 21, 22, 23, 24, 25, 26, 27, 28, 29, 30, 31, 32, 33, 34, 35, 36, 
  37, 38, 39, 40, 41, 42, 43, 44, 45, 46, 47, 48, 49, 50, 51, 52))
\end{verbatim}

También mediante sentencias \texttt{print}, se puede obtener información general:

\begin{verbatim}
  >>> print(graph)
  Graph named 'berlin52' with 52 nodes and 1378 edges
\end{verbatim}

La lista de nodos:

\begin{verbatim}
  print(graph.nodes)
\end{verbatim}

Aunque esta orden da el mismo resultado que \texttt{print(list(problem.get\_nodes()))}, en forma de lista.

%%% Local Variables:
%%% TeX-master: "../book"
%%% End:
